\documentclass[a4paper,12pt]{article}
\usepackage{amsmath}
\usepackage{graphicx}
\usepackage{geometry}
\geometry{top=1in,bottom=1in,left=1in,right=1in}

\title{Mini-Project \# \\}
\author{Group \# \\ Student Name(s) and SCIPER(s)\\}
\date{}

\begin{document}

\maketitle

\section*{1. Introduction}
Brief description of the problem statement.\\
Objectives of the project.\\
FEA Software Used: e.g. MATLAB\\

\section*{2. Methodology}
\subsection*{Model Setup}
Define material properties, geometry, and boundary conditions.\\
Analysis Type: (Modal, Transient, or Harmonic).\\
Assumptions and simplifications.

\section*{3. Results and Discussion}
Meshing, mode shapes, deformation, or stress distribution.\\
Tables and graphs of natural frequencies, time-history response, or frequency response.\\
Interpretation of the results.\\
Comparison with theoretical calculations or reference data.

\section*{4. Conclusion}
Summary of key findings.\\
Design improvements or future enhancements.

\section*{5. References}
\begin{itemize}
    \item Author(s). Title of the paper/book. \textit{Journal/Publisher}. (Year).
    \item Author(s). Title of the paper/book. \textit{Journal/Publisher}. (Year). 
\end{itemize}

\section*{Appendix (If required)}
Additional screenshots, detailed calculations, or raw data.

\end{document}
